\chapter{Bandzuordnung und Quellen}
\label{b21guide:ch:sources}

Der Leitfaden verwendet Begriffe und Notation aus Band~21
mit eigener Nummerierung. Die folgende Zuordnung nennt auch
Quellen für ergänzende Aussagen, die im Hauptband nicht stehen.

\begin{center}
\small
\begin{tabularx}{\linewidth}{@{}p{.37\linewidth}X@{}}
\toprule
Themenbereich & Zuordnung\\
\midrule
Endliche Folgen und Rekursion &
Band~21; natürliche Indizes und Rekursion aus Band~10.\\
Konvergenz, Abschätzungen, Rechenregeln &
Band~21; \cite[Abschnitte~2.1--2.2]{lebl};
\cite[Sätze~3.23, 3.26]{hunter}.\\
Teilfolgen, Häufungswerte, Cauchy-Kriterium &
Band~21; \cite[Abschnitte~2.3--2.4]{lebl};
\cite[Sätze~3.46, 3.57--3.58]{hunter}.\\
Monotone Konvergenz, bestimmte Divergenz &
Band~21; \cite[Satz~2.1.10 und Abschnitt~2.3]{lebl}.\\
Geometrische Folgen und Eulersche Folge &
\cite[Propositionen~2.2.11--2.2.12]{lebl};
\cite[Propositionen~3.31--3.32]{hunter}.\\
Cesàro-Mittel, Gewichtungen, Wurzelgrenzwerte &
\cite[Abschnitt~6.9]{semmes};
\cite[Druckseiten~1 und~3]{stolz}.\\
Geometrische Fehlerabschätzung &
\cite[Aufgabe~2.4.2]{lebl}; quantitative Folgenfassung.\\
Stolz--Cesàro &
\cite{stolz}; Nullfolgenfassung in \cite{mathtutor}.\\
Stabile Fehlerrekursion &
\cite[Druckseite~4, Lemma, Teil~(1)]{berinde}; affine Folgerung.\\
Fakultätswachstum &
\cite[Beispiel~1.1.4]{yan}; Eulersche Folge und Wurzelkriterien.\\
Umkehrkriterium für Cesàro-Mittel &
\cite[Druckseite~100, Satz~8.9]{mcmullen}; Folgenfassung.\\
Endliche Mittelungsschemata &
\cite[Abschnitt~1]{fan}; endliche reelle Dreieckszeilen.\\
Periodenstruktur &
\cite[Abschnitt~4.7.1]{grinberg}.\\
Endliche Differenzen und Newtondarstellung &
\cite[Abschnitt~5.3, (5.44)]{concrete}.\\
Subadditives und multiplikatives Wachstum &
\cite[PDF-Seite~10]{bruin}; Fekete und seine Varianten.\\
Endliche Summationsformeln &
\cite[Seite~1]{li}; verschobene Indizes und nichtnegative Gewichte.\\
Endliche monotone Teilfolgen &
\cite[Seiten~467--468]{erdos-szekeres}; scharfe Schranke.\\
\bottomrule
\end{tabularx}
\end{center}

Die Quellen werden in der Notation des Leitfadens verwendet;
weiterführende Strukturtheorie wird nicht übernommen.
Ergänzende Fassungen der Perioden-, Index-, Häufungswert- und
Mittelwertsätze sowie quantitative Rekursionsschranken und die
Fakultätswurzel sind Folgerungen der genannten Grundlagen.
Das gilt auch für bestimmte Divergenz und $a_n-\sigma_n\to0$
im tauberschen Kriterium sowie die superadditive und die
multiplikativen Fekete-Fassungen.
Diese Anpassungen sind keine wörtlichen Quellenaussagen.

\renewcommand{\bibname}{Literatur}
\begin{thebibliography}{99}
\addcontentsline{toc}{section}{Literatur}
\small
\setlength{\itemsep}{3pt}
\setlength{\parskip}{0pt}

\bibitem{project}
Martin Kunze: \emph{Die Grundlagen der Mathematik}.
Band~21: Folgen; Band~10: natürliche Zahlen und Rekursion;
Band~19: reelle Zahlen.

\bibitem{lebl}
Jiří Lebl: \emph{Basic Analysis I: Introduction to Real Analysis}.
Autorenfassung, Kapitel~2, Abschnitte~2.1--2.4.
\url{https://www.jirka.org/ra/html/sec_seqsandlims.html};
\url{https://www.jirka.org/ra/html/sec_factslimsseqs.html};
\url{https://www.jirka.org/ra/html/sec_bw.html};
\url{https://www.jirka.org/ra/html/sec_cauchy.html}.

\bibitem{hunter}
John K. Hunter: \emph{Introduction to Analysis}.
University of California, Davis, Kapitel~3: \emph{Sequences},
Druckseiten~35--59.
\url{https://www.math.ucdavis.edu/~hunter/intro_analysis_pdf/ch3.pdf}.

\bibitem{semmes}
Stephen Semmes: \emph{Some basic topics in analysis}.
Rice University, Abschnitt~6.9: \emph{Cesàro sums},
insbesondere Druckseite~121 zu arithmetischen Mitteln von Folgen.
\url{https://math.rice.edu/~semmes/math322.pdf}.

\bibitem{stolz}
Gabriel Nagy: \emph{The Stolz-Cesaro Theorem}.
Kansas State University, Druckseite~1: gewichtete Mittel;
Druckseite~2: allgemeine Fassung
des Stolz--Cesàro-Satzes, additiver Cesàro-Satz und
asymptotische Steigung; Druckseite~3: geometrische Mittel,
Quotienten- und Wurzelgrenzwerte.
\url{https://www.math.ksu.edu/~nagy/snippets/stolz-cesaro.pdf}.

\bibitem{mathtutor}
Math Tutor, Czech Technical University in Prague:
\emph{Sequences -- Theory -- Limit: Methods of evaluation}.
Abschnitt \emph{Stolz-Cesaro Theorem}, insbesondere die
Fassung für zwei gegen $0$ konvergente Folgen.
\url{https://math.fel.cvut.cz/mt/txta/2/txe3aa2h.htm}.


\bibitem{berinde}
Vasile Berinde und \c{S}tefan M\u{a}ru\c{s}ter:
\emph{Recurrent Inequalities and Some of their Applications to
Picard and Mann Iterations}. 2014, Druckseite~4:
Lemma zur Rekursionsungleichung~(1.1), Teil~(1).
\url{https://doi.org/10.2478/awutm-2014-0001}.

\bibitem{yan}
Min Yan: \emph{Advanced Analysis}.
Hong Kong University of Science and Technology,
Beispiel~1.1.4, Druckseite~10: $c^n/n!\to0$.
\url{https://www.math.hkust.edu.hk/~mamyan/ma4033/lecture.pdf}.

\newpage
\bibitem{mcmullen}
Curtis T. McMullen: \emph{Advanced Real Analysis}.
Harvard University, Math~212a, Druckseite~100, Satz~8.9
und vorangehende Aussage zu Cesàro-Konvergenz.
\url{https://people.math.harvard.edu/~ctm/home/text/class/harvard/212a/19/html/home/course/course.pdf}.

\bibitem{fan}
Steve Fan: \emph{Summability and the Closed Graph Theorem}.
Dartmouth College, Abschnitt~1, Druckseite~1:
Silverman--Toeplitz-Theorem.
\url{https://math.dartmouth.edu/~stevefan/papers/Summability%20and%20the%20Closed%20Graph%20Theorem.pdf}.


\bibitem{grinberg}
Darij Grinberg: \emph{Math 235: Mathematical Problem Solving}.
Autorenfassung, 20.~September~2026, Abschnitt~4.7.1,
Theorem~4.7.8 und Corollary~4.7.9, Druckseiten~167 und~169.
\url{https://darijgrinberg.gitlab.io/t/20f/mps.pdf}.

\bibitem{concrete}
Ronald L. Graham, Donald E. Knuth und Oren Patashnik:
\emph{Concrete Mathematics}, 2.~Auflage, Abschnitt~5.3,
Druckseiten~189--191, insbesondere Gleichung~(5.44).
\url{https://www.pearson.de/media/muster/ext/9780134389967.pdf}.

\bibitem{bruin}
Henk Bruin: \emph{Ergodic Theory I}, Vorlesung~12,
Universität Wien, PDF-Seite~10: \emph{Fekete's Lemma}.
\url{https://www.mat.univie.ac.at/~bruin/ET1_lect12.pdf}.


\bibitem{li}
Kin Y. Li: \emph{Summation by Parts}.
Mathematical Excalibur~11(3), 2006, Seite~1:
endliche Summationsformel und \emph{Abel's Inequality}.
\url{https://www.math.hkust.edu.hk/excalibur/v11_n3.pdf}.

\bibitem{erdos-szekeres}
Paul Erdős und George Szekeres:
\emph{A combinatorial problem in geometry}.
Compositio Mathematica~2 (1935), 463--470;
hier nur der Satz über monotone Teilfolgen und seine scharfe
Schranke auf den Seiten~467--468, insbesondere~(8).
\url{https://www.numdam.org/item/CM_1935__2__463_0.pdf}.

\end{thebibliography}
